\paragraph{模障碍的导子化闭包、共振角色群的同调分类与 Ihara--Chebotarev 误差字典：从 cyclotomic 因子到 \texorpdfstring{$p$}{p}-典型整性与零温线性响应}
在有限状态加权图/同步核的扭曲谱框架中，“模障碍/共振”并非散点经验现象，而可被压缩为一组有限维、整系数、可计算的同调—谱证书：同调导子 \(d\) 决定全部例外模数；共振角色构成有限阿贝尔群并可由 Smith 标准形一次性输出；非共振通道上存在显式相位失配下界给出统一的扭曲谱降；一旦共振发生，行列式必含 cyclotomic 因子并诱导有限阶旋转模态，从而在时间域出现不可消去的严格周期振荡校正。
与之并行，Ihara--Chebotarev 的平方根级误差与谱隙条件互为充要，而 Ihara--Mertens 常数与局部 \(p\)-因子均可经导数/同余闭包被压缩为可审计的有限维对象；在零温端，Perron 根的一阶微扰系数进一步给出一条与谱隙正交的“易激发性”标量序。
\begin{enumerate}
  \item \textbf{有限导子原理：全部模障碍由单一整数 \texorpdfstring{$d$}{d} 完全控制。}
  对原始有向图 \(G\) 的整数边权 \(w\)，回路权和集合 \(\mathcal{C}\) 的最大公因子
  \(d=\gcd(\mathcal{C})\)
  给出模共边界的充要判别：存在 \(g:V\to\ZZ/q\ZZ\) 使
  \(w(e)\equiv g(\mathrm{head}(e))-g(\mathrm{tail}(e))\ (\mathrm{mod}\ q)\)
  当且仅当 \(q\mid d\)（命题 \ref{prop:gm-cycle-gcd-mod-obstruction}）。
  因此所谓“无限多模数共振”必等价于一个固定的同调障碍 \(d>1\)，而非解析偶然。

  \item \textbf{统一扭曲谱降的定量下界：非障碍模上谱隙由最小相位失配显式给出。}
  若 \(q\nmid d\)，则扭曲邻接算子满足
  \[
  \boxed{\ \rho\!\bigl(A(\chi_q)\bigr)\ \le\ \bigl(1-\kappa_G\eta(q)\bigr)\rho\ },
  \qquad
  \eta(q)=\min_{1\le t<q}\sin^2\!\Bigl(\frac{\pi t}{q}\Bigr)\asymp q^{-2}
  \]
  其中 \(\kappa_G>0\) 仅依赖底图的原始性常数（定理 \ref{thm:gm-uniform-twist-gap-from-gcd}）。
  该不等式把“主弧衰减/谱压制”的来源压缩为有限维谱半径比较与显式三角函数下界的组合。

  \item \textbf{共振角色群的完整同调分类与可计算性：共振不是散点而是有限阿贝尔群。}
  共振角色集合
  \(\mathcal{X}_{\mathrm{res}}=\{\chi:\rho(A(\chi))=\rho(A(1))\}\)
  构成有限阿贝尔群，并同构于同调商挠性部分的 Pontryagin 对偶
  \[
  \boxed{\ \mathcal{X}_{\mathrm{res}}\ \cong\ \mathrm{Hom}\!\Bigl(\mathrm{tors}\bigl(H_1(G,\ZZ)/\langle[w]\rangle\bigr),\ \QQ/\ZZ\Bigr)\ }
  \]
  （定理 \ref{thm:gm-resonant-character-group-homology}）。
  其不变因子分解与生成角色表可由整数矩阵的 Smith 标准形在多项式时间内输出（命题 \ref{prop:gm-smith-compute-obstructions}），从而把“共振结构”提升为一次性可审计的离散输出。

  \item \textbf{cyclotomic 因子 \texorpdfstring{$\Longleftrightarrow$}{<->} 等半径共振：共振必然对应有限阶旋转模态并诱导严格周期振荡项。}
  在图 \(\zeta\)/Ihara--Artin \(\zeta\) 的行列式框架中，存在非平凡共振通道当且仅当扭曲 \(\zeta\)（或等价的 \(\det(I-uA(\chi))\)、\(\det(I-tB_\rho(u))\)）含 cyclotomic 因子；并且该情形等价于某个非平凡扭曲满足谱半径等半径
  \(\rho(A(\chi))=\rho(A(1))\)
  或
  \(\rho(B_\rho(u))=\lambda_1(u)\)
  （定理 \ref{thm:gm-graph-zeta-cyclotomic-classification}；定理 \ref{thm:conclusion75-ihara-cyclotomic-factor-classification}）。
  cyclotomic 因子对应单位圆上的单位根，从而把谱边界上的特征值强制为有限阶旋转；其时间域后果是：主项之外出现与 \(\lambda_1(u)^n/n\) 同阶的周期振荡校正项，其周期为相应单位根的阶。

  \item \textbf{统一扭曲谱隙的平凡共振判别与可计算闭包。}
  定义“统一扭曲谱隙”性质为：存在常数 \(\delta>0\) 使对一切 \(\chi\neq \mathbf 1\) 有
  \[
  \rho(A(\chi))\le (1-\delta)\rho(A(1)).
  \]
  则该性质当且仅当 \(\mathcal X_{\mathrm{res}}=\{\mathbf 1\}\)。
  事实上，若统一谱隙成立则不存在等半径共振角色；
  反之若 \(\mathcal X_{\mathrm{res}}=\{\mathbf 1\}\)，则由上一条的 cyclotomic 分类与障碍闭子群分解，非平凡扭曲被分为“障碍上非平凡”的无穷族与“障碍上平凡”的有限族：
  前者由定理 \ref{thm:gm-uniform-twist-gap-from-gcd} 给出统一谱降常数；
  后者由于有限且无等半径点，取最大谱半径比即可得到 \(\delta_1>0\)，合并得统一 \(\delta=\min(\delta_0,\delta_1)\)。
  由结论 \ref{con:xi-resonant-character-snf-computable} 的 SNF 可计算性，\(\mathcal X_{\mathrm{res}}\) 的平凡性可在多项式时间内审计，从而“统一扭曲谱隙是否成立”在该离散口径下属于确定性多项式时间可判定问题。

  \item \textbf{共振见证的素阶压缩：最小反证据必为素数阶。}
  若 \(\mathcal X_{\mathrm{res}}\neq\{\mathbf 1\}\)，则存在某个 \(\chi\in\mathcal X_{\mathrm{res}}\setminus\{\mathbf 1\}\) 其阶为素数。
  证明只需在有限阿贝尔群 \(\mathcal X_{\mathrm{res}}\) 中取阶最小的非平凡元：若其阶为合成数 \(m=ab\) 且 \(1<a<m\)，则 \(\chi^a\neq \mathbf 1\) 且 \(\mathrm{ord}(\chi^a)=b<m\)，与极小性矛盾，故该最小阶必为素数。
  因此一旦统一谱隙失败，总可在某个素阶通道上给出等半径共振的有限核验证据，从而把“失败见证”的搜索空间压缩到素数模数族。

  \item \textbf{Ihara--Chebotarev 的谱隙充要性：平方根级误差与 Ramanujan 型谱界互为等价判别。}
  对按 Frobenius 共轭类分层的 Ihara-primitive 计数 \(q_{n,C}(u)\)，平方根级一致误差
  \[
  q_{n,C}(u)=\frac{|C|}{|G|}\cdot\frac{\lambda_1(u)^n}{n}
  +O\!\left(\frac{\lambda_1(u)^{n/2}}{n}\right)
  \qquad(\text{对 }C\text{ 一致})
  \]
  与谱隙条件 \(\Lambda(u)\le \lambda_1(u)^{1/2}\) 互为充要（推论 \ref{cor:ihara-chebotarev-grh-criterion}）。
  因此误差尺度本身可作为谱隙质量的判别器，而无需先验知道全部扭曲谱。

  \item \textbf{Ihara--Mertens 常数的 Möbius--Abel 公式：主极点留数可由行列式导数给出有限维证书。}
  在 Perron--Frobenius 谱隙条件下，Ihara 口径的 Abel 有限部常数 \(\log\mathfrak{M}^{\mathrm{Ih}}_K(u)\) 具有与未回溯常数层同型的 Möbius--Abel 快速配方，并且主极点留数常数由
  \(
  \partial_t\det(I-tB(u))|_{t=t_\star(u)}
  \)
  的有限维导数证书完全决定（命题 \ref{prop:ihara-mertens-constant}）。

  \item \textbf{\texorpdfstring{$p$}{p}-典型 \texorpdfstring{$p$}{p}-进整性闭包：从 Dwork 同余到局部 \texorpdfstring{$p$}{p}-因子并自动稳定于多势参数。}
  Dwork 型同余与 Dieudonn\'e--Dwork 整性判据给出 \(p\)-典型局部因子
  \(
  Z_p(T;u)=\exp(\sum_{k\ge 0}a_{p^k}(u)T^{p^k}/p^k)
  \)
  在 \(1+T\ZZ_p[u][[T]]\) 内闭合（定理 \ref{thm:witt-dwork-zp-integrality}；引理 \ref{lem:witt-dieudonne-dwork}）。
  同时，多维势下的 necklace 整除与多维 Dwork 同余逐字成立（\S\ref{app:multi-dim-witt-guardrails}），因此“局部 \(p\)-因子”结构在系数环升级时不引入额外审计代价。

  \item \textbf{零温线性响应的代数刚性：一阶 carry 密度完全由 carry-free 骨架决定并导出可比标量序。}
  若 carry-free 骨架 \(A_0\) primitive 且 Perron 根 \(\lambda_0\) 简单，则谱半径在 \(u=0\) 处解析并满足
  \[
  \lambda_K(u)=\lambda_0+c_1u+O(u^2),
  \qquad
  c_1=\frac{\ell^{\mathsf T}A_1 r}{\ell^{\mathsf T}r},
  \qquad
  \mathbb{E}_u[\kappa]=\frac{c_1}{\lambda_0}u+O(u^2)
  \]
  （命题 \ref{prop:lambda-linear-response} 与推论 \ref{cor:kappa-linear-response-order}）。
  在 \(K_9/K_{13}/K_{21}\) 三核算例中，该一阶斜率给出零温附近的“抗 carry”排序，并且 \(K_{21}\) 的一阶响应常数在代数上塌缩为 \(c_1=\sqrt5\)（推论 \ref{cor:k21-c1-sqrt5}）。
\end{enumerate}

\begin{theorem}[carry 势零温极限的支撑刚性]\label{thm:conclusion-carry-zero-temp-support-collapse}
设 \(M_K(u)\) 为有限图上的 carry 加权转移矩阵，并取局部常势
\[
\varphi_\beta:=-\beta\kappa,
\qquad
\beta\to+\infty.
\]
令 \(\mu_\beta\) 为相应 Gibbs 平衡态。则任意零温极限点 \(\mu_0\) 都满足
\[
\boxed{\;\mu_0(\kappa=0)=1.\;}
\]
亦即，零温极限的全部质量必定落在 carry-free 子图生成的有限并 SFT 上。若 carry-free 子图仅有唯一的本质强连通分量，则 \(\mu_0\) 唯一，且正是该分量的 Parry 测度。
\end{theorem}
\begin{proof}
正文已建立谱半径导数与平衡态 carry 密度之间的恒等式
\[
\partial_\theta\log\lambda_K(e^\theta)=\mathbb E_{\mu_{e^\theta}}[\kappa].
\]
同时，关于零温极限的结论表明：当 \(\beta\to+\infty\) 时，任意极限平衡态都必须支撑在使势函数 \(-\kappa\) 取最大值的轨道上；而这恰等价于
\[
\kappa=0.
\]
故任意零温极限 \(\mu_0\) 都满足 \(\mu_0(\kappa=0)=1\)，从而其支撑落在 carry-free 子图所生成的有限并 SFT 上。

若 carry-free 子图只有一个本质强连通分量，则正文的零温选择定理进一步表明极限平衡态唯一，并等于该 primitive 分量的 Parry 最大熵测度。证毕。
\end{proof}

